\documentclass[11pt]{article}

% ---------------------------------------------------------------------------
% Self-contained: minimal package set, compiles with any modern TeX.
% ---------------------------------------------------------------------------
\usepackage[a4paper,margin=25mm]{geometry}
\usepackage{amsmath,amssymb}
\usepackage{booktabs}
\usepackage{array}
\usepackage[hidelinks]{hyperref}

\newcommand{\psib}{\psi}
\newcommand{\gradpsi}{\lVert\nabla\psi\rVert}
\newcommand{\code}[1]{\texttt{#1}}

\title{Failure of iterative defect(-error) correction (iDEC) for
gradient-injecting departure-point reconstructions in the
semi-Lagrangian level-set method}
\author{leia OpenFOAM module --- research record\\
\small (generated as a hand-off note for a follow-up DEC investigation)}
\date{2026-07-10}

\begin{document}
\maketitle

\begin{abstract}
We record, for a follow-up investigation, exactly \emph{what} failed with an
iterative defect(-error) correction (iDEC) reconstruction in the semi-Lagrangian
(SL) level-set solver \code{leiaSemiLagrangeLevelSetFoam}, \emph{for which cases},
and \emph{how}. The short version: a departure-point reconstruction that injects a
separately-computed cell gradient into the interpolation (\code{defectCorrectedIDW}
and its two iterative variants) is \textbf{unstable} under repeated SL transport on
all three verification flows (2D reversed shear vortex, 3D shear, 3D LeVeque
deformation, both $\mathrm{CFL}=0.5$ and $1.0$): the signed-distance gradient
$\gradpsi$ grows without bound while the interface shape error fails to converge.
Iterating the correction does \emph{not} fix it. We show numerically that (i) the
converged fixed point of the nodal-defect iteration still amplifies, and (ii) the
advection-level deferred correction (a Picard iteration on the reconstruction)
\emph{accelerates} the divergence. The diagnosis is that the reconstruction
operator has spectral radius $\rho>1$; iterating an amplifying operator cannot
stabilise it. The value-fitting \code{quadraticWeightedLeastSquares} (no injected
gradient) is stable and convergent on the same flows and serves as the control.
This document is intended to seed a separate effort on defect/error-correction
(DEC) formulations that could be made contractive.
\end{abstract}

\section{Context and problem statement}

The SL update advects the level set $\psi$ along characteristics,
\begin{equation}
\psi^{\,n+1}(\mathbf{x}_c) \;=\; \psi^{\,n}(\mathbf{x}_d),
\qquad
\mathbf{x}_d \;=\; \mathbf{x}_c - \mathbf{u}^{n+1}\Delta t
+ \tfrac12\!\left[\tfrac{\mathbf{u}^{n+1}-\mathbf{u}^{n}}{\Delta t}
+ (\mathbf{u}^{n+1}\!\cdot\nabla)\mathbf{u}^{n+1}\right]\!\Delta t^2 ,
\end{equation}
where $\mathbf{x}_d$ is the second-order Taylor backward-foot (fixed; not the
subject here) and the reconstruction of $\psi^{\,n}$ at $\mathbf{x}_d$ is
runtime-selectable (\code{slReconstruction}). Every reconstruction must reproduce
the arrival-cell value, $\mathcal{R}(c,\mathbf{x}_c)=\psi^{\,n}_c$, and is built on
the arrival cell's cell-point-cell (CPC) point-neighbour stencil
$S(c)=\{c\}\cup\{\text{point neighbours}\}$ (${\sim}9$ in 2D, ${\sim}27$ in 3D).
No reinitialisation is performed, by design, so $\gradpsi=1$ is \emph{not} enforced;
it is expected to drift but must stay \emph{bounded}.

Errors reported below (aggregated per resolution $h=1/N$):
\emph{shape} $=$ geometric interface (marker-traced) error at the reversal time $T$;
\emph{gradient} $=\lVert\gradpsi-1\rVert_{L_2}$ in the narrow band;
\emph{volume} $=$ phase-volume (mass) conservation error.

\section{The reconstruction and its two iterative variants}

\subsection{Single-pass gradient-enhanced Shepard (\code{defectCorrectedIDW})}
Inverse-distance-weighted (IDW) interpolation lifted to second order by a per-node
Taylor/defect term with the cell-centred least-squares gradient
$\mathbf{g}_C=\nabla\psi^{\,n}_C$ (OpenFOAM \code{pointCellsLeastSquares}):
\begin{equation}
\psi(\mathbf{x}_d) \;=\;
\frac{\sum_{C\in S} w_C\,\big[\psi_C + \mathbf{g}_C\!\cdot(\mathbf{x}_d-\mathbf{x}_C)\big]}
     {\sum_{C\in S} w_C},
\qquad w_C=\lvert \mathbf{x}_d-\mathbf{x}_C\rvert^{-p}.
\label{eq:singlepass}
\end{equation}
As a \emph{static} interpolant \eqref{eq:singlepass} is linear-exact (unit test:
error $\sim 3\times10^{-16}$). It is the composition of a bounded Shepard smoother
with an \emph{independently computed} differentiator $\mathbf{g}_C$.

\subsection{Variant A: iterative nodal-defect correction}
Derived (via an independent-derivation panel) as the ``correct'' iterative reading:
solve, per cell, for effective nodal values $u_C=\psi_C+\delta_C$ so that a
\emph{regularised}, row-stochastic blend $A$ reproduces the stencil data,
\begin{equation}
\sum_{C} A_{DC}\,\big(u_C + \mathbf{g}_C\!\cdot(\mathbf{x}_D-\mathbf{x}_C)\big)=\psi_D
\quad\forall D\in S,
\qquad
A_{DC}=\frac{w^{\mathrm{reg}}(\mathbf{x}_D-\mathbf{x}_C)}{\sum_E w^{\mathrm{reg}}(\mathbf{x}_D-\mathbf{x}_E)},
\end{equation}
by damped Gauss--Seidel (dict controls \code{dcIters}$=3$, \code{dcRelax}$=0.8$,
\code{dcRegEps}$=0.1$; $A_{DD}>\tfrac12$ makes the \emph{inner} solve a contraction),
with the arrival-node defect pinned $\delta_0=0$ (preserves the contract). The final
foot value uses the original singular weights on the corrected Taylor values. The
regularisation is essential: with singular weights $A=I$ and the nodal defect
collapses to zero, recovering \eqref{eq:singlepass}.

\subsection{Variant B: advection-level deferred correction}
A Picard/deferred correction wrapped around the reconstruction in
\code{slAdvection::advect()} (the backtracking foot is computed once and reused,
byte-identical): iterate
\begin{equation}
\psi^{(0)}=\psi^{\,n},\qquad
\psi^{(k+1)}(\mathbf{x}_c)=
\mathcal{R}\!\big[\psi^{(k)}\big](\mathbf{x}_d(c)),
\label{eq:picard}
\end{equation}
so the gradient $\mathbf{g}$ is recomputed from the current iterate each pass
(``self-consistent gradient''). Controls \code{nDefCorr}, \code{defCorrRelax},
\code{defCorrTol}; \code{nDefCorr}$=1$ reproduces the single pass exactly.

\section{Numerical results: the failure}

All runs: reversed (oscillating) flows so $\text{exact}(T)=\text{initial}$;
detrixheAslam phase indicator; $N=32/64/128$; the reported columns are at the
finest available resolution and $\mathrm{CFL}=0.5$ unless noted.

\subsection{Single pass \eqref{eq:singlepass}: $\gradpsi$ blows up}
\begin{center}
\begin{tabular}{l r r r}
\toprule
case & shape($T$) trend & gradient $\lVert\gradpsi-1\rVert_{L_2}$ (as $h\downarrow$) & shape order \\
\midrule
2D reversed vortex & $1.6\!\times\!10^{-2}\!\to\!3.0\!\times\!10^{-2}$ & $39 \to 1.4\!\times\!10^{4}\to 4.7\!\times\!10^{9}$ & $-0.25$ \\
3D shear           & $0.83 \to 0.92$ (interface destroyed) & $8.5\!\times\!10^{3}\to 1.2\!\times\!10^{9}$ & $\approx 0$ \\
3D deformation     & $0.45 \to 0.48$ (interface destroyed) & $8.8\!\times\!10^{4}\to 4.5\!\times\!10^{12}$ & $\approx 0$ \\
\bottomrule
\end{tabular}
\end{center}
The always-on value cap in the advection keeps $\psi$ \emph{values} bounded (so the
shape error saturates rather than NaN-ing), but $\gradpsi$ sawtooths to enormous
values --- the signed-distance field is destroyed.

\subsection{Variant A (iterative nodal-defect): helps, but still diverges}
\begin{center}
\begin{tabular}{l r r r}
\toprule
case & shape($T$) $N{=}32\!\to\!128$ & gradient $N{=}32\!\to\!128$ & shape order \\
\midrule
2D reversed vortex & $1.16\!\times\!10^{-2}\!\to\!2.55\!\times\!10^{-2}$ & $15.2 \to 1.66\!\times\!10^{3}\to 6.3\!\times\!10^{7}$ & $-0.33$ \\
3D shear           & $0.341 \to 0.610$ ($N{=}64$) & $50.1 \to 1.15\!\times\!10^{4}$ ($N{=}64$) & $-0.84$ \\
3D deformation     & $0.302 \to 0.399$ ($N{=}64$) & $2.36\!\times\!10^{2}\to 3.39\!\times\!10^{5}$ ($N{=}64$) & $-0.99$ \\
\bottomrule
\end{tabular}
\end{center}
Iterating reduces the amplification constant by $1$--$2$ orders (2D $N{=}128$:
$6.3\!\times\!10^{7}$ vs.\ $4.7\!\times\!10^{9}$ single-pass) but the character is
unchanged: the error \emph{grows} under refinement.

\paragraph{The iteration is converged, not under-iterated.}
Same case (2D, $N{=}64$, $\mathrm{CFL}=0.5$, $\lVert\gradpsi-1\rVert_{L_2}$ at $T$):
\begin{center}
\begin{tabular}{l r}
\toprule
Variant A setting & gradient error \\
\midrule
\code{dcIters}$=3$ (default)      & $1657$ \\
\code{dcIters}$=30$               & $1649$ \quad(\emph{identical}: GS at its fixed point in 3 sweeps) \\
\code{limitSlope}\,=\,true        & $1657$ \quad(Barth--Jespersen increment limiter: no effect) \\
\bottomrule
\end{tabular}
\end{center}
More sweeps change nothing --- the \emph{fixed point itself} amplifies --- and the
slope limiter (which bounds the increment at neighbour centres) does not bound the
grid-scale $\gradpsi$ oscillation.

\subsection{Variant B (advection-level deferred correction): diverges faster}
Same case, sweeping the number of deferred-correction passes:
\begin{center}
\begin{tabular}{l r}
\toprule
\code{nDefCorr} & gradient error $\lVert\gradpsi-1\rVert_{L_2}$ at $T$ \\
\midrule
$1$ (single pass)             & $1.66\!\times\!10^{3}$ \\
$3$                           & $4.33\!\times\!10^{12}$ \\
$5$                           & $2.93\!\times\!10^{22}$ \\
$10$                          & $1.25\!\times\!10^{47}$ \\
$10$, \code{defCorrRelax}$=0.5$ & $5.62\!\times\!10^{18}$ \\
\bottomrule
\end{tabular}
\end{center}
Each pass multiplies the error by roughly the operator gain: \code{nDefCorr}$=K$
blows up ${\sim}K$ times faster. Under-relaxation slows but does not cross below the
divergence threshold.

\subsection{Control: value-fitting \code{quadraticWeightedLeastSquares} is stable}
Same flows, same solver, \emph{no injected gradient} (weighted least-squares
\emph{value} fit; its slope is a by-product, never an input):
\begin{center}
\begin{tabular}{l r r r}
\toprule
case & shape($T$) $N{=}32\!\to\!128$ & shape order (LSQ) & gradient (bounded) \\
\midrule
2D reversed vortex & $3.0\!\times\!10^{-4}\!\to\!4.7\!\times\!10^{-6}$ & $3.01$ & ${\sim}0.1$--$0.2$ \\
3D shear           & $1.5\!\times\!10^{-3}\!\to\!5.2\!\times\!10^{-5}$ & $2.43$ & ${\sim}0.38$ \\
3D deformation     & $5.3\!\times\!10^{-3}\!\to\!8.2\!\times\!10^{-4}$ & $1.34$ & ${\sim}0.4$--$0.5$ \\
\bottomrule
\end{tabular}
\end{center}
(The lower 3D-deformation order is under-resolution of the thin LeVeque filament,
not instability: the per-interval order \emph{rises} with refinement and the volume
error converges $\sim$2nd order.)

\section{Diagnosis: the operator has spectral radius $\rho>1$}
The SL step is a fixed-point map $\psi^{\,n+1}=\mathcal{M}[\psi^{\,n}]$. For the
gradient-injecting reconstructions \eqref{eq:singlepass}, $\mathcal{M}$ is
(bounded Shepard smoother) $\circ$ (\emph{unbounded} least-squares differentiator).
Differentiation amplifies grid-scale ($\sim$ Nyquist) content; the smoother does not
remove it. The lift term $\mathbf{g}_C\!\cdot(\mathbf{x}_d-\mathbf{x}_C)$ feeds the
gradient of $\psi^{\,n}$ back into $\psi^{\,n+1}$ without ever being constrained by
the reconstructed values, so $\lVert\mathcal{M}'\rVert>1$ on the high-frequency
modes: $\rho(\mathcal{M})>1$ and $\gradpsi$ diverges geometrically in the step index
(the measured $10^{9}$--$10^{12}$ over one reversal period). This is the same failure
family as \code{linearTaylor} and \code{nestedLSQ}, which also inject an LSQ gradient.

\section{Why iteration cannot help}
A defect/deferred correction stabilises only a \emph{contracting} base operator
($\rho\le1$): it drives the iterate to the fixed point of a stable map.
\begin{itemize}
\item \textbf{Variant A} drives the \emph{nodal} defect to zero (its \emph{inner}
Gauss--Seidel is a contraction, $A_{DD}>\tfrac12$), but the resulting reconstruction
is still evaluated with the singular-weight blend of $\psi_C+\mathbf{g}_C\!\cdot r$
--- the gradient path is not severed --- so the \emph{outer} SL map still has
$\rho>1$. Hence the converged fixed point amplifies and more sweeps do nothing.
\item \textbf{Variant B} iterates the SL map itself \eqref{eq:picard}. Iterating a
map with $\rho>1$ is Picard on a divergent fixed point: it \emph{amplifies faster}
(the measured $K$-fold acceleration).
\end{itemize}
In both cases the lower-order-gradient ``lift'' does not ``only work iteratively'';
the object being iterated is expansive, not contractive.

\section{Directions for a DEC investigation}
The instability is intrinsic to injecting a \emph{value-inconsistent}
externally-differentiated gradient. Promising directions to make a defect/error
correction contractive (each untested here):
\begin{enumerate}
\item \textbf{Consistent-gradient DEC.} Replace the external LSQ gradient by the
\emph{gradient of the reconstructed polynomial itself} (as \code{quadraticWeightedLeastSquares}
does implicitly), so the lift cannot inject an independent mode. This is essentially
why the value fit is stable; a DEC that keeps value and slope self-consistent per
iterate may inherit $\rho\le1$.
\item \textbf{Bounded / limited gradient.} Constrain $\mathbf{g}_C\!\cdot r$ to the
stencil range (a genuine slope limiter on the \emph{lift}, not on the increment at
neighbour centres as the current \code{limitSlope}), trading order for a guaranteed
$\rho\le1$; then defect-correct \emph{upward} in order from the limited, stable base.
\item \textbf{Low-order stable base $\oplus$ deferred high-order defect
(Khosla--Rubin form).} Use plain IDW/donor (monotone, $\rho\le1$) as the operator
that is \emph{inverted} each pass, with the high-order term as a lagged
right-hand-side defect --- but only if the base is genuinely contracting (plain IDW
in the SL loop must be verified to be).
\item \textbf{BFECC / error-compensated SL.} Round-trip (advect forward then
backward), estimate the transport error, and compensate --- a defect correction on
the \emph{trajectory/advection} rather than on the reconstruction; second-order and
known to be stable for SL when the base advection is monotone.
\item \textbf{Discrete-exterior-calculus (DEC) discretisation.} If ``DEC'' is meant
as discrete exterior calculus, a structure-preserving gradient/codifferential on the
dual mesh may give a better-conditioned (non-amplifying) gradient operator than the
cell LSQ; worth checking whether it changes $\rho$.
\end{enumerate}

\paragraph{Reproduction.} All numbers above come from the leia SL suite
(\code{leiaSemiLagrangeLevelSetFoam}) driven by the snakemake workflow. The three
diagnostic sweeps (\code{dcIters}, \code{limitSlope}, \code{nDefCorr}) were run on
the 2D reversed vortex at $N{=}64$, $\mathrm{CFL}=0.5$; the per-case tables from the
\code{idwCompare\{2D,3Dshear,3Ddeformation\}} studies. The reconstruction models and
controls live in \code{src/leiaLevelSet/semiLagrangian/}
(\code{defectCorrectedIDWReconstruction}, \code{slCorrector}/\code{deferredCorrector},
\code{quadraticWeightedLeastSquaresReconstruction}); the Taylor backtracking in
\code{slAdvection.C} was never modified.

\end{document}
